\documentclass[11pt]{article}
\usepackage[margin=1in]{geometry}
\usepackage{amsmath,amssymb,amsfonts,amsthm,mathtools}
\usepackage{booktabs,array}
\usepackage[hidelinks]{hyperref}
\usepackage{enumitem}
\usepackage{microtype}

\title{GQE Reward Design}
\author{}
\date{}

\begin{document}
\maketitle

\section{Goal}

We consider quantum unitary synthesis as a multi-objective optimization problem with three competing goals:
\begin{enumerate}[label=(\roman*)]
    \item maximize process fidelity,
    \item minimize circuit depth,
    \item minimize CNOT count.
\end{enumerate}

The central challenge is that these objectives should \emph{not} be collapsed too aggressively into a naive weighted sum. In particular:
\begin{itemize}
    \item fidelity is the primary objective,
    \item structural simplification should matter more once fidelity is already good,
    \item small temporary degradations in fidelity may occasionally be useful for exploring new regions of the Pareto front,
    \item the reward should encourage the discovery of \emph{new non-dominated tradeoffs}, not only improvements relative to a single incumbent.
\end{itemize}

This document proposes a reward with two parts:
\begin{enumerate}[label=(\arabic*)]
    \item a \textbf{local comparison reward}, which compares a candidate circuit against a reference circuit and provides a dense learning signal, and
    \item a \textbf{Pareto-front bonus}, which rewards expansion of the global non-dominated archive.
\end{enumerate}

The full reward is therefore both \emph{locally shaped} and \emph{globally Pareto-aware}.

\section{Notation}

Let the candidate circuit be denoted by $x$, and let the reference circuit be denoted by $x_{\mathrm{ref}}$.
For each circuit we define:
\begin{align}
    F &\in [0,1] &&\text{candidate process fidelity,} \\
    D &\in \{0,1,2,\dots\} &&\text{candidate circuit depth,} \\
    C &\in \{0,1,2,\dots\} &&\text{candidate CNOT count,}
\end{align}
and similarly
\begin{align}
    F_{\mathrm{ref}},\quad D_{\mathrm{ref}},\quad C_{\mathrm{ref}}
\end{align}
for the reference circuit.

We assume that the optimization direction is:
\begin{itemize}
    \item maximize $F$,
    \item minimize $D$,
    \item minimize $C$.
\end{itemize}

We also introduce the following small positive constants:
\begin{align}
    \varepsilon_{\phi} &> 0 &&\text{stabilizer for log-infidelity,} \\
    \delta_{\mathrm{bad}} &\ge 0 &&\text{tolerance defining a clearly bad fidelity drop,} \\
    \tau &> 0 &&\text{temperature of the soft fidelity gate.}
\end{align}

\section{Why Raw Fidelity Differences Are Not Enough}

Using the raw fidelity difference $F-F_{\mathrm{ref}}$ is often too crude for unitary synthesis.
For example, improving fidelity from $0.99$ to $0.999$ is usually much more valuable than improving from $0.50$ to $0.509$, even though both changes have similar absolute magnitude.

To reflect this, we work with \emph{log-infidelity utility}.
Define
\begin{equation}
    \Phi(F) = -\log(1-F+\varepsilon_{\phi}).
\end{equation}

This function has the following properties:
\begin{enumerate}[label=(\alph*)]
    \item it is increasing in $F$,
    \item it grows rapidly as $F \to 1$,
    \item it transforms improvements near perfect fidelity into larger reward increments.
\end{enumerate}

We then define the fidelity improvement relative to the reference as
\begin{equation}
    \Delta_{\Phi}(x\mid x_{\mathrm{ref}})
    = \Phi(F) - \Phi(F_{\mathrm{ref}}).
\end{equation}

Interpretation:
\begin{itemize}
    \item $\Delta_{\Phi} > 0$: the candidate improves fidelity relative to the reference,
    \item $\Delta_{\Phi} = 0$: the candidate matches the reference in transformed fidelity,
    \item $\Delta_{\Phi} < 0$: the candidate is worse in fidelity.
\end{itemize}

\section{Structure Improvement Terms}

We want structural improvements to be measured in a scale-aware way. A reduction from depth $100$ to $90$ should not be treated exactly the same as a reduction from $10$ to $0$.
A natural choice is to use log-ratios.

\subsection{Positive structure gains}

Define the positive depth and CNOT improvements by
\begin{align}
    G_D(x\mid x_{\mathrm{ref}})
    &= \left[\log\frac{D_{\mathrm{ref}}+1}{D+1}\right]_+, \\
    G_C(x\mid x_{\mathrm{ref}})
    &= \left[\log\frac{C_{\mathrm{ref}}+1}{C+1}\right]_+,
\end{align}
where
\begin{equation}
    [z]_+ = \max(z,0).
\end{equation}

Thus:
\begin{itemize}
    \item $G_D > 0$ only if $D < D_{\mathrm{ref}}$,
    \item $G_C > 0$ only if $C < C_{\mathrm{ref}}$,
    \item both gains are zero when the corresponding structure metric does not improve.
\end{itemize}

\subsection{Optional structure regression penalties}

It is often useful to also penalize unnecessary structural regressions whenever fidelity is not worse.
Define
\begin{align}
    P_D(x\mid x_{\mathrm{ref}})
    &= \left[\log\frac{D+1}{D_{\mathrm{ref}}+1}\right]_+, \\
    P_C(x\mid x_{\mathrm{ref}})
    &= \left[\log\frac{C+1}{C_{\mathrm{ref}}+1}\right]_+.
\end{align}

Thus:
\begin{itemize}
    \item $P_D > 0$ only if $D > D_{\mathrm{ref}}$,
    \item $P_C > 0$ only if $C > C_{\mathrm{ref}}$.
\end{itemize}

These terms are optional. If one wants a softer fidelity-first reward, they can be disabled by setting the corresponding coefficients to zero.

\section{Soft Fidelity Gate}

We do not want structure to matter equally in all regimes.
If the candidate has very low fidelity, rewarding depth or CNOT reduction too strongly may push the search toward shallow but useless circuits.
On the other hand, once fidelity is already good, structural simplification becomes much more meaningful.

To model this smoothly, define the fidelity gate
\begin{equation}
    w(F) = \sigma\!\left(\frac{F-F_{\mathrm{gate}}}{\tau}\right),
\end{equation}
where
\begin{equation}
    \sigma(z)=\frac{1}{1+e^{-z}}
\end{equation}
is the logistic sigmoid, $F_{\mathrm{gate}}\in(0,1)$ is a chosen fidelity threshold, and $\tau>0$ controls the sharpness.

Interpretation:
\begin{itemize}
    \item if $F \ll F_{\mathrm{gate}}$, then $w(F)\approx 0$ and structure bonuses are strongly suppressed,
    \item if $F \gg F_{\mathrm{gate}}$, then $w(F)\approx 1$ and structure bonuses are almost fully active,
    \item for intermediate fidelities, structure pressure turns on gradually rather than abruptly.
\end{itemize}

This is preferable to a hard threshold because it avoids discontinuous reward jumps.

\section{Local Comparison Reward}

The local reward compares the candidate $x$ to a reference $x_{\mathrm{ref}}$.
It is piecewise defined so that:
\begin{enumerate}[label=(\roman*)]
    \item large fidelity drops are penalized strongly,
    \item small fidelity drops are still discouraged, but not so harshly that exploration dies,
    \item non-worse fidelity activates stronger structure optimization.
\end{enumerate}

\subsection{Definition}

Let
\begin{align}
    \alpha_{\mathrm{bad}},\alpha_{\mathrm{soft}},\alpha_{\mathrm{good}} &> 0, \\
    \kappa_{\mathrm{bad}},\kappa_{\mathrm{soft}},\kappa_{\mathrm{good}} &> 0, \\
    \lambda_D,\lambda_C &\ge 0, \\
    \mu_D,\mu_C &\ge 0, \\
    \eta &\in [0,1].
\end{align}

Then define
\begin{equation}
R_{\mathrm{local}}(x\mid x_{\mathrm{ref}})=
\begin{cases}
-\alpha_{\mathrm{bad}}
\left(
    e^{\kappa_{\mathrm{bad}}(-\Delta_{\Phi}-\delta_{\mathrm{bad}})}-1
\right),
& \Delta_{\Phi}< -\delta_{\mathrm{bad}}, \\
-\alpha_{\mathrm{soft}}
\left(
    e^{\kappa_{\mathrm{soft}}(-\Delta_{\Phi})}-1
\right)
+
\eta\,w(F_{\mathrm{ref}})
\left(
    e^{\lambda_D G_D + \lambda_C G_C}-1
\right),
& -\delta_{\mathrm{bad}}\le \Delta_{\Phi}<0, \\
\alpha_{\mathrm{good}}
\left(
    e^{\kappa_{\mathrm{good}}\Delta_{\Phi}}-1
\right)
+
 w(F)
\left(
    e^{\lambda_D G_D + \lambda_C G_C}-1
\right)
-
 w(F)(\mu_D P_D + \mu_C P_C),
& \Delta_{\Phi}\ge 0.
\end{cases}
\label{eq:Rlocal}
\end{equation}

\subsection{Meaning of every term}

We now explain every term in \eqref{eq:Rlocal}.

\paragraph{Regime 1: clearly bad fidelity drop.}
If
\begin{equation}
    \Delta_{\Phi}< -\delta_{\mathrm{bad}},
\end{equation}
then the reward is
\begin{equation}
    -\alpha_{\mathrm{bad}}\Bigl(e^{\kappa_{\mathrm{bad}}(-\Delta_{\Phi}-\delta_{\mathrm{bad}})}-1\Bigr).
\end{equation}
Here:
\begin{itemize}
    \item $-\Delta_{\Phi}-\delta_{\mathrm{bad}}>0$ measures how far below the tolerated fidelity region the candidate falls,
    \item $\kappa_{\mathrm{bad}}$ controls how sharply the penalty grows,
    \item $\alpha_{\mathrm{bad}}$ sets the overall magnitude of this penalty.
\end{itemize}
This branch gives a \emph{strictly negative} reward and ignores structure entirely.

\paragraph{Regime 2: small fidelity drop.}
If
\begin{equation}
    -\delta_{\mathrm{bad}}\le \Delta_{\Phi}<0,
\end{equation}
then the reward is the sum of:
\begin{enumerate}[label=(\alph*)]
    \item a soft penalty
    \begin{equation}
        -\alpha_{\mathrm{soft}}\Bigl(e^{\kappa_{\mathrm{soft}}(-\Delta_{\Phi})}-1\Bigr),
    \end{equation}
    which discourages any fidelity loss, and
    \item a reduced structure bonus
    \begin{equation}
        \eta\,w(F_{\mathrm{ref}})\Bigl(e^{\lambda_D G_D+\lambda_C G_C}-1\Bigr),
    \end{equation}
    where $\eta\in[0,1]$ attenuates structure reward in this exploratory regime.
\end{enumerate}

This branch captures the idea that slight fidelity degradation may sometimes be tolerated for exploration, but it should remain costly overall unless the structural gain is genuinely promising.

\paragraph{Regime 3: non-worse fidelity.}
If
\begin{equation}
    \Delta_{\Phi}\ge 0,
\end{equation}
then the reward has three components.

First, the fidelity-improvement reward:
\begin{equation}
    \alpha_{\mathrm{good}}\Bigl(e^{\kappa_{\mathrm{good}}\Delta_{\Phi}}-1\Bigr).
\end{equation}
This is zero when fidelity is unchanged and positive when fidelity improves. Its exponential form makes high-end fidelity gains increasingly valuable.

Second, the structure-improvement bonus:
\begin{equation}
    w(F)\Bigl(e^{\lambda_D G_D+\lambda_C G_C}-1\Bigr).
\end{equation}
This is gated by current fidelity through $w(F)$, so structure matters more once the candidate is already in a high-fidelity regime.

Third, the optional structure-regression penalty:
\begin{equation}
    -w(F)(\mu_D P_D + \mu_C P_C).
\end{equation}
This discourages unnecessary depth or CNOT increases when fidelity is not worse.

\section{Why Exponential Coupling Is Useful}

The term
\begin{equation}
    e^{\lambda_D G_D+\lambda_C G_C}-1
\end{equation}
implements a nonlinear coupling between structural improvements.
It has several useful properties:
\begin{enumerate}[label=(\arabic*)]
    \item it is zero when there is no structural improvement,
    \item it is increasing in both $G_D$ and $G_C$,
    \item simultaneous improvements in depth and CNOT receive more reward than either alone,
    \item small improvements remain gentle, while stronger improvements receive superlinear emphasis.
\end{enumerate}

This is in the same spirit as nonlinear scalarizations used in architecture search, where one does not want a simple linear sum to flatten the tradeoff geometry too aggressively.

\section{Pareto-Front Bonus}

The local reward alone still compares the candidate only to a single reference circuit.
However, true multi-objective optimization should also reward \emph{global frontier expansion}.

Let $A_t$ denote the archive of non-dominated circuits before adding the current candidate.
We define a Pareto-front bonus by
\begin{equation}
    R_{\mathrm{pareto}}(x;A_t)=\beta_{\mathrm{HV}}\,\Delta HV(x;A_t),
\label{eq:RparetoHV}
\end{equation}
where $\beta_{\mathrm{HV}}\ge 0$ is a scaling coefficient and $\Delta HV(x;A_t)$ is the increase in 3-dimensional hypervolume obtained by adding $x$ to the archive.

The three objectives are:
\begin{itemize}
    \item maximize fidelity,
    \item minimize depth,
    \item minimize CNOT count.
\end{itemize}

Thus:
\begin{itemize}
    \item if $x$ is dominated by the archive, then typically $\Delta HV(x;A_t)=0$,
    \item if $x$ is a new useful tradeoff point, then $\Delta HV(x;A_t)>0$,
    \item the more the candidate expands the frontier, the larger the bonus.
\end{itemize}

\subsection{Cheaper Pareto proxy}

If exact 3-dimensional hypervolume is inconvenient to compute at every step, one can use a simpler proxy:
\begin{equation}
    R_{\mathrm{pareto\text{-}proxy}}(x;A_t)
    =
    \beta_1\,\mathbf{1}[x\text{ is non-dominated w.r.t. }A_t]
    +
    \beta_2\,N_{\mathrm{dom}}(x;A_t),
\label{eq:RparetoProxy}
\end{equation}
where:
\begin{itemize}
    \item $\beta_1\ge 0$ rewards the candidate for entering the frontier,
    \item $\beta_2\ge 0$ rewards the candidate for dominating existing archive points,
    \item $N_{\mathrm{dom}}(x;A_t)$ is the number of archive points dominated by $x$.
\end{itemize}

This proxy is cruder than hypervolume improvement, but it is easy to implement and often works well in practice.

\section{Total Reward and Cost}

The full reward is
\begin{equation}
    R_{\mathrm{total}}(x)
    =
    R_{\mathrm{local}}(x\mid x_{\mathrm{ref}})
    +
    R_{\mathrm{pareto}}(x;A_t),
\label{eq:Rtotal}
\end{equation}
or, if using the cheaper archive proxy,
\begin{equation}
    R_{\mathrm{total}}(x)
    =
    R_{\mathrm{local}}(x\mid x_{\mathrm{ref}})
    +
    R_{\mathrm{pareto\text{-}proxy}}(x;A_t).
\end{equation}

If the training code is written in terms of a cost to minimize, we simply define
\begin{equation}
    \mathrm{cost}(x) = -R_{\mathrm{total}}(x).
\label{eq:cost}
\end{equation}

Thus maximizing reward and minimizing cost are equivalent.

\section{Behavior Analysis}

We now analyze how the reward behaves in important scenarios.

\subsection{Scenario 1: fidelity improves, structure unchanged}

Suppose
\begin{equation}
    \Delta_{\Phi}>0,
    \qquad
    D=D_{\mathrm{ref}},
    \qquad
    C=C_{\mathrm{ref}}.
\end{equation}
Then
\begin{equation}
    G_D=G_C=P_D=P_C=0,
\end{equation}
and the local reward reduces to a strictly positive fidelity term:
\begin{equation}
    R_{\mathrm{local}}=
    \alpha_{\mathrm{good}}\Bigl(e^{\kappa_{\mathrm{good}}\Delta_{\Phi}}-1\Bigr)>0.
\end{equation}
So fidelity improvement is always rewarded.

\subsection{Scenario 2: fidelity preserved, structure improves}

Suppose
\begin{equation}
    \Delta_{\Phi}=0,
    \qquad
    D<D_{\mathrm{ref}}
    \quad\text{and/or}\quad
    C<C_{\mathrm{ref}}.
\end{equation}
Then the fidelity term is zero, but the structure term is positive:
\begin{equation}
    R_{\mathrm{local}}=w(F)\Bigl(e^{\lambda_D G_D+\lambda_C G_C}-1\Bigr)>0.
\end{equation}
Hence structurally simpler circuits are preferred among equal-fidelity candidates.

\subsection{Scenario 3: large fidelity drop, even with better structure}

Suppose
\begin{equation}
    \Delta_{\Phi}< -\delta_{\mathrm{bad}}.
\end{equation}
Then
\begin{equation}
    R_{\mathrm{local}}<0
\end{equation}
regardless of $D$ and $C$.
This preserves the fidelity-first principle in the region judged to be clearly unsafe.

\subsection{Scenario 4: small fidelity drop but major structural gain}

Suppose
\begin{equation}
    -\delta_{\mathrm{bad}}\le \Delta_{\Phi}<0,
\end{equation}
and assume that $D$ and/or $C$ improve substantially.
Then the reward becomes
\begin{equation}
    R_{\mathrm{local}}=
    \underbrace{-\alpha_{\mathrm{soft}}\Bigl(e^{\kappa_{\mathrm{soft}}(-\Delta_{\Phi})}-1\Bigr)}_{\text{negative}}
    +
    \underbrace{\eta w(F_{\mathrm{ref}})\Bigl(e^{\lambda_D G_D+\lambda_C G_C}-1\Bigr)}_{\text{positive}}.
\end{equation}
This regime allows controlled exploration.
Depending on hyperparameter choices, the total may remain negative, near zero, or slightly positive.
That flexibility is intentional: it allows the search to occasionally pass through simpler intermediate circuits that may later lead to better high-fidelity solutions.

\subsection{Scenario 5: fidelity improves but structure worsens}

Suppose
\begin{equation}
    \Delta_{\Phi}>0,
    \qquad
    D>D_{\mathrm{ref}}
    \quad\text{or}\quad
    C>C_{\mathrm{ref}}.
\end{equation}
Then the candidate still receives a positive fidelity reward, but may be penalized through $P_D$ and $P_C$:
\begin{equation}
    R_{\mathrm{local}}=
    \alpha_{\mathrm{good}}\Bigl(e^{\kappa_{\mathrm{good}}\Delta_{\Phi}}-1\Bigr)
    -w(F)(\mu_D P_D+\mu_C P_C)
    + \text{possible structure gains}.
\end{equation}
This means fidelity is still primary, but unnecessary structural regressions are not ignored.

\subsection{Scenario 6: candidate is a new frontier point}

Even if the local reward is only modestly positive, the candidate can receive extra reward through
\begin{equation}
    R_{\mathrm{pareto}}(x;A_t)>0
\end{equation}
when it expands the non-dominated archive.
This is one of the most important properties of the full design: it does not only reward improvement against a single incumbent, but also rewards the discovery of novel tradeoff points.

\section{Edge Cases}

\subsection{Edge case: $F=0$}

If $F=0$, then
\begin{equation}
    \Phi(0) = -\log(1+\varepsilon_{\phi}-1) = -\log(1+\varepsilon_{\phi}-1?)
\end{equation}
which is more transparently written as
\begin{equation}
    \Phi(0) = -\log(1-0+\varepsilon_{\phi}) = -\log(1+\varepsilon_{\phi}).
\end{equation}
Since $\varepsilon_{\phi}>0$ is small, this is finite and slightly negative. There is no singularity.

If the reference has much better fidelity, then $\Delta_{\Phi}$ will be strongly negative and the candidate falls into the bad-drop regime.

\subsection{Edge case: $F\to 1$}

As $F\to 1$,
\begin{equation}
    \Phi(F) = -\log(1-F+\varepsilon_{\phi})
\end{equation}
becomes large but remains finite because of $\varepsilon_{\phi}$.
This is desirable: very high fidelity improvements are emphasized without causing numerical blow-up.

\subsection{Edge case: $D=0$ or $C=0$}

Because all depth and CNOT ratios are written using $D+1$ and $C+1$, there is no division by zero:
\begin{align}
    G_D &= \left[\log\frac{D_{\mathrm{ref}}+1}{D+1}\right]_+, \\
    G_C &= \left[\log\frac{C_{\mathrm{ref}}+1}{C+1}\right]_+.
\end{align}
So empty or trivial circuits are handled safely.

\subsection{Edge case: candidate identical to reference}

If
\begin{equation}
    F=F_{\mathrm{ref}},
    \qquad
    D=D_{\mathrm{ref}},
    \qquad
    C=C_{\mathrm{ref}},
\end{equation}
then
\begin{equation}
    \Delta_{\Phi}=0,
    \quad
    G_D=G_C=P_D=P_C=0.
\end{equation}
Hence
\begin{equation}
    R_{\mathrm{local}}=0.
\end{equation}
If the archive bonus is also zero, then the total reward is exactly zero. This is desirable: no change gives no reward.

\subsection{Edge case: fidelity equal, depth better, CNOT worse}

Suppose
\begin{equation}
    \Delta_{\Phi}=0,
    \qquad
    D<D_{\mathrm{ref}},
    \qquad
    C>C_{\mathrm{ref}}.
\end{equation}
Then both a positive gain term $G_D$ and a regression penalty $P_C$ may be active.
The reward will trade these off according to $\lambda_D$ and $\mu_C$.
This is appropriate because the candidate is a mixed tradeoff point rather than a uniformly better circuit.

\subsection{Edge case: structure improves strongly at very low fidelity}

Suppose the candidate is structurally much simpler but fidelity is poor. Then $w(F)$ is close to zero, which suppresses structure bonuses. This prevents the optimization from overvaluing shallow but inaccurate circuits.

\section{Hyperparameter Interpretation}

The main hyperparameters play the following roles:
\begin{center}
\begin{tabular}{>{\raggedright\arraybackslash}p{0.18\linewidth}>{\raggedright\arraybackslash}p{0.72\linewidth}}
\toprule
Parameter & Meaning \\
\midrule
$\varepsilon_{\phi}$ & numerical stabilizer in log-infidelity; should be small \\
$\delta_{\mathrm{bad}}$ & tolerance separating mild from clearly bad fidelity drops \\
$F_{\mathrm{gate}}$ & fidelity level at which structure reward begins to matter strongly \\
$\tau$ & softness of the fidelity gate; smaller values make it closer to a hard threshold \\
$\alpha_{\mathrm{bad}}$ & scale of severe fidelity penalties \\
$\kappa_{\mathrm{bad}}$ & steepness of severe fidelity penalties \\
$\alpha_{\mathrm{soft}}$ & scale of mild fidelity penalties \\
$\kappa_{\mathrm{soft}}$ & steepness of mild fidelity penalties \\
$\alpha_{\mathrm{good}}$ & scale of fidelity-improvement reward \\
$\kappa_{\mathrm{good}}$ & steepness of fidelity-improvement reward \\
$\lambda_D,\lambda_C$ & strengths of structure-improvement bonuses \\
$\mu_D,\mu_C$ & strengths of structure-regression penalties \\
$\eta$ & amount of structure reward allowed during mild-fidelity-drop exploration \\
$\beta_{\mathrm{HV}}$ & weight of global Pareto expansion bonus \\
\bottomrule
\end{tabular}
\end{center}

\section{Recommended Qualitative Tuning Strategy}

A sensible practical strategy is:
\begin{enumerate}[label=(\arabic*)]
    \item start with strong fidelity protection, i.e. choose $\alpha_{\mathrm{bad}}$ and $\kappa_{\mathrm{bad}}$ so that clearly bad fidelity drops are always unattractive,
    \item make the mild-drop regime narrow by using a small $\delta_{\mathrm{bad}}$,
    \item set $\eta$ small so that exploration through mild degradation is possible but limited,
    \item choose $F_{\mathrm{gate}}$ fairly high so that depth and CNOT pressure turn on mostly in the medium-to-high fidelity regime,
    \item use a nonzero archive bonus so that novel Pareto points are not ignored.
\end{enumerate}

\section{What This Reward Guarantees and What It Does Not}

\subsection{What it guarantees}

The proposed reward guarantees the following qualitative behavior:
\begin{enumerate}[label=(\roman*)]
    \item large fidelity drops are punished strongly,
    \item structure matters more when fidelity is already good,
    \item equal-fidelity solutions are compared meaningfully by depth and CNOT count,
    \item the search can still discover useful non-dominated tradeoffs through archive-aware reward,
    \item the objective is more faithful to Pareto-front optimization than a simple weighted sum.
\end{enumerate}

\subsection{What it intentionally allows}

The reward intentionally allows a narrow regime in which a small fidelity drop can coexist with some structure reward. This is not a flaw. It is a deliberate exploration mechanism motivated by the fact that some structurally simpler intermediate circuits may later lead to better solutions.

\subsection{What it does not fully solve}

No scalar reward can perfectly represent the full Pareto structure of a multi-objective problem. The archive-aware bonus helps substantially, but final model selection should still inspect the actual non-dominated set rather than relying on the scalar reward alone.

\section{Bottom Line}

The proposed reward is
\begin{equation*}
    R_{\mathrm{total}}(x)
    =
    R_{\mathrm{local}}(x\mid x_{\mathrm{ref}})
    +
    R_{\mathrm{pareto}}(x;A_t),
\end{equation*}
with $R_{\mathrm{local}}$ defined in \eqref{eq:Rlocal} and $R_{\mathrm{pareto}}$ given either by hypervolume improvement \eqref{eq:RparetoHV} or by the cheaper proxy \eqref{eq:RparetoProxy}.

This design is a good fit for multi-objective unitary synthesis because it combines:
\begin{itemize}
    \item \textbf{fidelity sensitivity} through log-infidelity utility,
    \item \textbf{nonlinear structure reward} through exponential coupling,
    \item \textbf{controlled exploration} through a mild-drop regime,
    \item \textbf{Pareto awareness} through archive expansion bonus.
\end{itemize}

In short, it is not merely a fidelity-first reward and not merely a scalarized compression reward. It is a compromise between safe local improvement and true frontier discovery.

\end{document}
